% M2_NnaemekaNnachiEgwu.tex
% Milestone 2: Scientific Contextualization and Tradeoff Analysis
% Real-Time Systems Seminar, HSHL, Prof. Dr. Stefan Henkler
% Author: Nnaemeka Nnachi-Egwu
%
% Structured notes format per seminar checklist. Not publication prose.

\documentclass[conference,onecolumn]{IEEEtran}
\usepackage{amsmath, amssymb, booktabs, array, microtype, url}

\begin{document}

\title{Milestone 2: Scientific Contextualization\\
  {\large Priority Inheritance Protocol (Resource Access Protocols)}}

\author{\IEEEauthorblockN{Nnaemeka Nnachi-Egwu}
\IEEEauthorblockA{Electronic Engineering\\
Hochschule Hamm-Lippstadt\\
Lippstadt, Germany\\
nnaemeka.nnachi-egwu@stud.hshl.de}}

\maketitle

%---------------------------------------------------------------
\section{Protocol Landscape and Chronology}

\textbf{Chronological development} (Buttazzo \S7.1 overview):
\begin{enumerate}
  \item \textbf{NPP} (Non-Preemptive Protocol): classical; no cited inventor.
        Task runs non-preemptively inside CS. Simple, over-blocks.
  \item \textbf{HLP/IPC} (Highest Locker Priority): classical.
        Raises priority at entry time to ceiling of entered semaphore.
        Requires ceiling declarations in code.
  \item \textbf{PIP} + \textbf{PCP}: Sha, Rajkumar, Lehoczky [SRL90],
        IEEE Trans.\ Computers 39(9):1175--1185, 1990.
        \emph{Same paper, two protocols}: PIP (§7.6) and PCP (§7.7)
        as a design progression.
  \item \textbf{SRP}: Baker [Bak91], Real-Time Systems 3(1):67--99, 1991.
        Extends PCP: multi-unit resources, EDF, stack sharing.
\end{enumerate}

\textbf{Conceptual positions:}
\begin{itemize}
  \item NPP/HLP block at \emph{activation time} (before CS entry). Pessimistic.
  \item PIP defers blocking to \emph{access time} (at \texttt{pi\_wait}).
        No blocking unless actual contention. Less pessimistic.
  \item PCP adds lock-grant rule on top of PIP:
        prevents deadlock and chained blocking. Non-transparent.
  \item SRP blocks at \emph{preemption time} (like NPP/HLP) but with
        tighter dynamic ceilings; allows EDF and stack sharing.
\end{itemize}

%---------------------------------------------------------------
\section{Where PIP Sits Relative to Alternatives}

\textbf{PIP vs.\ NPP:}
NPP blocks all higher-priority tasks whenever \emph{any} CS is entered,
even when those tasks share no resource with the CS task.
PIP blocks only tasks that actually contend for the same semaphore.
More selective, less pessimistic.

\textbf{PIP vs.\ HLP:}
HLP blocks a task at \emph{activation time} based on ceiling of semaphore
being entered --- even if the branch accessing the resource is never taken
at runtime (conditional-CS pessimism, Buttazzo p.~214).
PIP defers to actual \texttt{wait} call. Removes conditional-branch pessimism.
However, HLP gives single blocking (Theorem~7.1) vs.\ PIP's $\alpha_i$
blockings (Theorem~7.2) --- HLP can be \emph{less} blocking in worst case.

\textbf{PIP vs.\ PCP:}

{\small\begin{tabular}{@{}lcc@{}}
\toprule
Property & PIP & PCP \\
\midrule
Deadlock-free & \textbf{No} & Yes (Thm~7.3) \\
Chained blocking & \textbf{Possible} & No (Thm~7.4) \\
Transparent & Yes & \textbf{No} \\
Max blockings & $\alpha_i = \min(l_i,s_i)$ & 1 \\
Blocking time & Tighter (actual contention) & Pessimistic (ceiling rule) \\
Impl.\ effort & Hard & Medium \\
\bottomrule
\end{tabular}}

\textbf{PIP vs.\ SRP (Baker [Bak91]):}
SRP extends PCP with: (1) multi-unit resources, (2) EDF support via
preemption levels, (3) runtime stack sharing (up to 90\% savings,
Buttazzo \S7.8.5). PIP is fixed-priority, single-unit, no stack sharing.

%---------------------------------------------------------------
\section{Kopetz: Priority Inversion as Correctness Failure}

\textbf{Kopetz definition (RTS slide 7):}
``A real-time computer system is one in which correctness depends
not only on logical results but also on the \emph{physical instant}
at which they are produced.''

\textbf{Implication:}
In the Fig.~7.4 scenario, $\tau_1$ has highest priority because its
deadline is tightest. If $\tau_1$ misses its deadline due to unbounded
inversion, the physical instant of its output is wrong --- not just late.
This is a system failure by definition, not a QoS degradation.

\textbf{Slide 13 (Prof.\ Henkler):} ``Real-time means predictable;
result after deadline is wrong, not just late.''

\textbf{Slide 48:} EDF optimal \emph{only without resource competition.}
PIP (or equivalent) is necessary to restore theoretical scheduling
guarantees in the presence of shared resources.

%---------------------------------------------------------------
\section{Application Domains}

\textbf{POSIX (\texttt{PTHREAD\_PRIO\_INHERIT}):}
\begin{itemize}
  \item POSIX.1-2008 standardises \texttt{pthread\_mutexattr\_setprotocol}
        with \texttt{PTHREAD\_PRIO\_INHERIT}.
  \item One attribute change at mutex creation; zero changes to lock/unlock sites.
  \item Supported: Linux (with RT patch / kernel $\geq$2.6.25), VxWorks, QNX.
  \item \textbf{Mars Pathfinder (1997):} VxWorks mutex without priority
        inheritance caused IPC bus watchdog resets. Fix: enable
        \texttt{PTHREAD\_PRIO\_INHERIT} on the meteorological bus mutex.
        One config change, no code rewrite.
\end{itemize}

\textbf{Automotive OSEK/AUTOSAR:}
\begin{itemize}
  \item OSEK OS default: ceiling protocol (HLP/IPC variant), not PIP.
  \item Some AUTOSAR vendor extensions and POSIX-compliant RTOS layers
        provide PIP-like inheritance (\textbf{qualified claim}: not a
        universal AUTOSAR requirement).
  \item ISO~26262 (ASIL): bounded blocking time required for safety
        qualification. PIP's $B_i$ bound + RTA provides this.
\end{itemize}

%---------------------------------------------------------------
\section{Where PIP Is Insufficient}

\textbf{Multiprocessor systems:}
\begin{itemize}
  \item Ch.~7 explicitly uniprocessor (\S7.3, p.~209).
  \item Boosting a holder's priority on a different processor requires
        inter-processor interrupt (IPI) with non-deterministic latency.
  \item Gracioli et~al.\ [GP] measure IPI overhead ($\Delta_\mathrm{ipi}$),
        cache affinity loss, scheduling overhead on actual hardware ---
        none appear in Buttazzo's $B_i$ formula.
  \item \textbf{Consequence:} PIP semantics cannot be directly applied
        on multiprocessors; MPCP / FMLP / MSRP needed.
\end{itemize}

\textbf{Distributed systems:}
Resource access becomes message-passing; priority inheritance requires
network communication; timing is non-deterministic. Kopetz [Ko] discusses
distributed timing constraints; PIP provides no help here.

\textbf{EDF scheduling:}
PIP defined for fixed priorities. For EDF systems, SRP [Bak91]
is the standard protocol (handles dynamic priorities via preemption levels).

%---------------------------------------------------------------
\section{Scientific and Engineering Tradeoffs}

{\small\begin{tabular}{@{}lp{5cm}p{4.5cm}@{}}
\toprule
Tradeoff & PIP advantage & PIP disadvantage \\
\midrule
Pessimism vs.\ correctness & Lower pessimism: blocks only on genuine contention & Multiple blockings ($\alpha_i$) vs.\ single blocking in HLP/PCP \\
Transparency vs.\ safety & No code change needed; legacy-friendly & Dynamic inheritance hard to audit statically; limits formal certification \\
Simplicity vs.\ completeness & No deadlock risk from protocol itself & Programmer must enforce total semaphore ordering \\
Generality vs.\ scope & Works for any fixed-priority RTOS & Uniprocessor only; no EDF support \\
\bottomrule
\end{tabular}}

%---------------------------------------------------------------
\section{Reference Justifications}

\textbf{Buttazzo 2011:} Primary source. Provides all formal
definitions, lemmas, theorems, algorithms, and the five-protocol
comparison.

\textbf{Sha, Rajkumar, Lehoczky 1990:} Original PIP/PCP paper
[SRL90]. Required to cite original authors of Theorem~7.2 and
PCP results.

\textbf{Kopetz 2011:} Grounding definition for real-time correctness
(physical instant, slide 7). Establishes why inversion is a
correctness failure, not a performance issue.

\textbf{Abella et~al.\ 2015:} WCET estimation limits. Contextualises
the assumption that $\delta_{i,k}$ values are exactly known.
None of the surveyed WCET methods is fully trustworthy.

\textbf{Gracioli et~al.\ 2013:} Multiprocessor OS overhead. Grounds
PIP's scope limitation with measured data from real hardware.

\textbf{Baker 1991:} SRP origin [Bak91]. Required for protocol
landscape (Table~7.5 SRP row) and distinguishing EDF-capable protocols.

\textbf{Behrmann et~al.\ 2004:} UPPAAL tutorial. Needed for M4
UPPAAL implementation and verification methodology.

\textbf{Excluded references (with justification):}
\begin{itemize}
  \item General scheduling textbooks (not PIP-specific; Buttazzo covers scope).
  \item AUTOSAR specifications (not publicly available; claim qualified instead).
  \item POSIX.1-2008 standard (too broad; PTHREAD\_PRIO\_INHERIT is a
        well-known fact, not requiring a primary citation in a 5-page paper).
\end{itemize}

\end{document}
